\documentclass[11pt]{article}

\usepackage[a4paper,margin=1in]{geometry}
\usepackage{amsmath,amssymb,amsthm,mathtools}
\usepackage{booktabs}
\usepackage{enumitem}
\usepackage{microtype}
\usepackage{xcolor}
\usepackage{hyperref}

\hypersetup{
    colorlinks=true,
    linkcolor=blue!60!black,
    citecolor=blue!60!black,
    urlcolor=blue!60!black,
    pdftitle={Factors Before Quotients},
    pdfauthor={Vinicius Santos}
}

\newtheorem{definition}{Definition}[section]
\newtheorem{theorem}[definition]{Theorem}
\newtheorem{proposition}[definition]{Proposition}
\newtheorem{corollary}[definition]{Corollary}
\newtheorem{lemma}[definition]{Lemma}
\newtheorem{remark}[definition]{Remark}
\newtheorem{question}[definition]{Question}

\newcommand{\B}{\mathsf{B}}
\newcommand{\Tr}{\operatorname{Tr}}
\newcommand{\Nm}{\operatorname{N}}
\newcommand{\Sp}{\operatorname{Sp}}
\newcommand{\Diff}{\operatorname{Diff}}
\newcommand{\sig}{\sigma}
\newcommand{\ph}{\varphi}
\newcommand{\ps}{\psi}
\newcommand{\Z}{\mathbb{Z}}
\newcommand{\Q}{\mathbb{Q}}
\newcommand{\R}{\mathbb{R}}
\newcommand{\C}{\mathbb{C}}
\newcommand{\dd}{\partial_{\ph}}
\newcommand{\Del}{\Delta_{\ph}}
\newcommand{\Cph}{C_{\ph}}
\newcommand{\Eph}{E_{\ph}}
\newcommand{\Etw}{\mathcal E_{\ph}}
\newcommand{\ord}{\operatorname{ord}}

\title{Factors Before Quotients\\
\large Golden Difference Calculus, Bounded Derivative Weights, and Two Exponential Layers}
\author{Vinicius Santos\\
\small Exploratory draft}
\date{June 2026}

\begin{document}
\maketitle

\begin{abstract}
The companion manuscripts \emph{A Golden Shadow Algebra}, \emph{Quotients Without Division}, \emph{Shadow Trigonometry}, and \emph{Signatures Before Roots} suggest a recurring method: operations that appear primitive in classical notation may become projections, relations, localizations, or lift problems in a $\ph$-native algebraic presentation.  Addition becomes trace projection.  Division splits into unit action, rate localization, relation, approximation, and factor projection.  Trigonometry begins from norm-one phase.  Roots appear as spread and signature lifts.

This note applies the same method to calculus.  Classically, differentiation is introduced as a quotient.  Here the primitive datum is the golden contraction
\[
    \Cph(x)=\ph^{-1}x=(\ph-1)x,
\]
which is already a native unit action generated by the scalar core.  For polynomials over $R=\Z[\ph]$, the golden difference
\[
    \Del f(x)=f(x)-f(\ph^{-1}x)
\]
is divisible by the golden displacement
\[
    \Del x=x-\ph^{-1}x.
\]
The golden derivative is therefore defined as the factor projection
\[
    \Del f=(\dd f)\Del x,
\]
not as a primitive pointwise quotient.

The resulting operator is formally the Jackson $q$-difference operator at $q=\ph^{-1}$.  The new point is not historical invention of $q$-calculus, but the $\ph$-native route to it and the golden arithmetic forced by this specific value of $q$.  The derivative weights satisfy
\[
    [n]_{\ph}=1+\ph^{-1}+\cdots+\ph^{-(n-1)}
       =\ph^2-\ph^{2-n}.
\]
Under the chosen real embedding $\ph=(1+\sqrt5)/2$, these weights are bounded and converge to $\ph^2$.  Under the conjugate embedding they are unbounded; boundedness is a property of the selected real place, not of the abstract ring.

This bounded-weight phenomenon has a direct analytic consequence: the untwisted eigen-equation $\dd E=E$ gives a reciprocal-product exponential with radius of convergence exactly $\ph^2=\ph+1$, hence a non-entire function.  By contrast, the twist-compatible equation $\dd\mathcal E=\sig_{\ph}(\mathcal E)$ gives a product exponential whose finite approximants are denominator-free polynomials in $\Z[\ph][x]$ and whose completed product is entire.  These are the two standard $q$-exponentials at $q=\ph^{-1}$; the contribution here is their placement in the $\ph$-native hierarchy.  The untwisted exponential lives in the localization layer, while the twisted exponential lives first in the polynomial/completion layer.  The same orbit viewpoint explains the logarithmic boundary: on the golden unit group, logarithm is the exact integer index $\ord_{\ph}(\ph^k)=k$; off the group, scalar analytic logarithm is an interpolation of that index, with denominators recording the cost of leaving the discrete orbit.  The product exponential has zeros at $-\ph^{2+k}$, the reciprocal exponential has poles at $+\ph^{2+k}$, and the nearest orbit obstruction explains the common radius $\ph^2$.  The integration layer has a parallel arithmetic obstruction: among the weights $[n]_{\ph}$, only $[1]_{\ph}$ and $[2]_{\ph}$ are units in $\Z[\ph]$.
\end{abstract}

\tableofcontents

\section{Introduction: calculus begins with divisibility}

The previous notes in this sequence repeatedly relocate classical operations into more structural layers \cite{santosGoldenShadow,santosQuotients,santosShadowTrig,santosSignatures}.  In \emph{A Golden Shadow Algebra}, the scalar grammar generated from the golden leaf by
\[
    \B(x,y)=xy-y
\]
was shown to be multiplication/predecessor-native but not addition-native.  Addition appeared only after passing to a typed pair layer:
\[
    x+y=\Tr\langle x\mid y\rangle.
\]
In \emph{Quotients Without Division}, division split into unit action, rate localization, quotient relation, reciprocal approximation, normalization constraint, and factor projection.  In \emph{Shadow Trigonometry}, sine and cosine were delayed in favor of unnormalized phase projections:
\[
    T(U)=U+U^{-1},
    \qquad
    S(U)=U-U^{-1}.
\]
In \emph{Signatures Before Roots}, square roots appeared as spread lifts from discriminant data:
\[
    \Sp^2=\Tr^2-4\Nm.
\]

Calculus is the next overloaded object.  The classical derivative is usually introduced as a quotient
\[
    \frac{f(x+h)-f(x)}{h},
\]
followed by a limiting process.  This notation bundles at least four distinct ideas: displacement, difference, divisibility by the displacement, and limiting scale.  The present note separates them.

The $\ph$-native displacement comes from the contraction already generated by the scalar core:
\[
    x\longmapsto \ph^{-1}x.
\]
For polynomials, the difference
\[
    f(x)-f(\ph^{-1}x)
\]
is divisible by
\[
    x-\ph^{-1}x.
\]
The derivative is the factor obtained from this divisibility relation.

The central slogan is:
\begin{quote}
Calculus begins not with division, but with divisibility.
\end{quote}

There is also a $\ph$-specific analytic payoff.  The monomial weights do not grow like the ordinary weights $n$.  Instead they collapse to a bounded golden orbit:
\[
    [n]_{\ph}=\ph^2-\ph^{2-n}\longrightarrow \ph^2
\]
in the chosen real embedding.  The same boundedness gives, by the ratio test, a non-entire reciprocal exponential with radius $\ph^2$.  But the twisted product rule also reveals a second, more polynomial-facing exponential: the twist-compatible eigen-equation produces denominator-free finite approximants and an entire completed product.  The two exponentials sort naturally into two layers rather than competing for a single definition.

\subsection*{Main results}

The main results are:
\begin{enumerate}[label=(\roman*)]
    \item For every polynomial $f\in \Z[\ph][x]$, there is a unique polynomial $\dd f\in\Z[\ph][x]$ satisfying
    \[
        f(x)-f(\ph^{-1}x)=(\dd f)(x)(x-\ph^{-1}x).
    \]
    \item The monomial rule is
    \[
        \dd x^n=[n]_{\ph}x^{n-1},
    \]
    where
    \[
        [n]_{\ph}=1+\ph^{-1}+\cdots+\ph^{-(n-1)}.
    \]
    \item The golden weights have the closed form
    \[
        [n]_{\ph}=\ph^2-\ph^{2-n}.
    \]
    In the chosen real embedding they increase to $\ph^2$; in the conjugate embedding they are unbounded.
    \item The untwisted eigen-equation $\dd E=E$ gives a reciprocal-product exponential with radius of convergence $\ph^2$, hence a non-entire function.
    \item The twisted eigen-equation $\dd\mathcal E=\sig_{\ph}(\mathcal E)$ gives a product exponential
    \[
        \mathcal E_{\ph}(x)=\prod_{k\ge0}(1+\ph^{-2-k}x)
    \]
    whose finite approximants are denominator-free polynomials and whose completed product is entire.
    \item The logarithmic structure collapses to one index seen in two regimes: on the golden unit group, $\ord_{\ph}(\ph^k)=k$ is exact and division-free; off the unit group, scalar analytic logarithm interpolates that index.  The poles of the reciprocal exponential occur at $+\ph^{2+k}$, while the zeros of the product exponential occur at $-\ph^{2+k}$.
    \item The product rule is the standard twisted $q$-Leibniz rule, with twist given by the native contraction $f(x)\mapsto f(\ph^{-1}x)$.
    \item The chain rule is not the ordinary chain rule.  It is a two-stage factorization, and the clean classical-looking formula holds for polynomial inner maps only in the contraction-compatible case $g(x)=cx$.
    \item Integration requires localization at the golden weights $[n]_{\ph}$; only $[1]_{\ph}$ and $[2]_{\ph}$ are units in $\Z[\ph]$.
\end{enumerate}

\section{The native hierarchy ladder}

This paper uses the layers developed in the preceding notes.  The point is not to replace all classical operations by one primitive, but to place each role in its proper layer.

\begin{center}
\small
\begin{tabular}{p{0.25\linewidth}p{0.30\linewidth}p{0.34\linewidth}}
\toprule
Layer & Role & Native mechanism \\
\midrule
Scalar layer & constants and unit actions & $R=\Z[\ph]$, $\B(x,y)=xy-y$ \\
Trace layer & addition & $x+y=\Tr\langle x\mid y\rangle$ \\
Rate layer & fixed denominators & classes $[x:s]_S\in S^{-1}R$ \\
Factor layer & quotient-like operators & extract factors from divisibility relations \\
Phase layer & trigonometry & norm-one phase, trace/spread projections \\
Signature layer & roots & spread lifts and discriminant signatures \\
Completion layer & infinite processes & limits, products, analytic functions \\
\bottomrule
\end{tabular}
\end{center}

The present note adds a calculus example to the factor layer.  Derivatives are quotients only after coordinate evaluation.  Algebraically, they are factor projections.  The exponential discussion then splits into two higher layers: localization for the reciprocal exponential, and completion for the product exponential.

\section{Golden contraction as unit action}

Let
\[
    q=\ph^{-1}=\ph-1.
\]
The scalar core already generates multiplication by $q$ because
\[
    \B(\ph,x)=x(\ph-1)=qx.
\]

\begin{definition}[Golden contraction]
The golden contraction is the unit action
\[
    \Cph(x)=qx=\ph^{-1}x.
\]
\end{definition}

This is not general division.  It is multiplication by a unit in $R=\Z[\ph]$.  The displacement of the coordinate $x$ under golden contraction is
\[
    \Del x=x-\Cph(x)=x-qx=(1-q)x.
\]
Since
\[
    1-q=1-\ph^{-1}=\ph^{-2},
\]
we have
\[
    \Del x=\ph^{-2}x.
\]
Thus even the displacement of $x$ is generated by native unit action.

\begin{remark}[Fixed scale]
The scale $q=\ph^{-1}$ is fixed.  This is not an infinitesimal calculus and it is not a $q\to1$ limit.  It is a self-similar finite-scale calculus adapted to the contraction already present in the $\ph$-scalar grammar.  Ordinary differential calculus appears only after a separate deformation in which $q\to1$.
\end{remark}

\section{Golden difference}

\begin{definition}[Golden difference]
For a polynomial $f\in R[x]$, define
\[
    \Del f(x)=f(x)-f(qx).
\]
\end{definition}

This is the difference of $f$ along the contraction orbit
\[
    x\mapsto qx.
\]
It is finite-scale and self-similar rather than infinitesimal.  The next section shows that for polynomials the difference is always divisible by the coordinate displacement $\Del x$.

\section{Derivative as factor projection}

\begin{definition}[Golden derivative]
For $f\in R[x]$, the golden derivative $\dd f$ is the polynomial factor, when it exists, satisfying
\[
    \Del f=(\dd f)\Del x.
\]
Equivalently,
\[
    f(x)-f(qx)=(\dd f)(x)(x-qx).
\]
\end{definition}

\begin{theorem}[Existence and uniqueness of the factor derivative]
Let $R=\Z[\ph]$ and $q=\ph^{-1}$.  For every $f\in R[x]$, there is a unique polynomial $\dd f\in R[x]$ such that
\[
    f(x)-f(qx)=(\dd f)(x)(x-qx).
\]
\end{theorem}

\begin{proof}
The polynomial $f(x)-f(qx)$ vanishes at $x=0$, hence is divisible by $x$ in $R[x]$.  Also
\[
    x-qx=(1-q)x=\ph^{-2}x,
\]
and $\ph^{-2}$ is a unit in $R$.  Therefore divisibility by $x$ is equivalent to divisibility by $x-qx$.  Existence follows.

For uniqueness, $R[x]$ is an integral domain and $x-qx\ne0$ because $q\ne1$.  If two polynomials $g,h$ satisfy
\[
    g(x)(x-qx)=h(x)(x-qx),
\]
then $(g-h)(x-qx)=0$, hence $g=h$.
\end{proof}

\begin{remark}[The fixed point at zero]
Pointwise, the classical quotient notation would have a $0/0$ appearance at $x=0$.  The factor projection avoids this coordinate singularity.  The derivative is a polynomial factor in $R[x]$, globally defined even at the fixed point of the contraction.  This is the same structural move used elsewhere in the series: the algebraic object is defined before coordinate evaluation makes it look singular.
\end{remark}

\begin{remark}[Why fixed scale matters]
If $q=1$, then $x-qx=0$ and the factor projection is undefined.  The construction works precisely because the golden scale is fixed away from $1$.  The same fact later explains why the reciprocal golden exponential is non-entire: the calculus is self-similar at one finite contraction scale, not infinitesimal.
\end{remark}

\section{The monomial rule}

For $f(x)=x^n$,
\[
    \Del x^n=x^n-(qx)^n=(1-q^n)x^n.
\]
Since
\[
    \Del x=(1-q)x,
\]
we have the factorization
\[
    1-q^n=(1-q)(1+q+\cdots+q^{n-1}).
\]

\begin{definition}[Golden integer weight]
For $n\ge1$, define
\[
    [n]_{\ph}=1+q+q^2+\cdots+q^{n-1},
    \qquad q=\ph^{-1}.
\]
Set $[0]_{\ph}=0$.
\end{definition}

\begin{proposition}[Monomial rule]
For all $n\ge1$,
\[
    \dd x^n=[n]_{\ph}x^{n-1}.
\]
\end{proposition}

\begin{proof}
Using the factorization above,
\[
    \Del x^n=(1-q^n)x^n
    =(1+q+\cdots+q^{n-1})x^{n-1}(1-q)x.
\]
Since $(1-q)x=\Del x$, the claimed factor is $[n]_{\ph}x^{n-1}$.
\end{proof}

This is the Jackson $q$-weight at $q=\ph^{-1}$, but the next section records the special collapse that occurs at the golden value.

\section{The golden collapse of derivative weights}

The formula
\[
    [n]_q=\frac{1-q^n}{1-q}
\]
is the generic $q$-integer.  The golden value $q=\ph^{-1}$ gives more structure because
\[
    1-\ph^{-1}=\ph^{-2}.
\]

\begin{proposition}[Closed form and bounded real-place weights]
For all $n\ge1$,
\[
    [n]_{\ph}=\ph^2-\ph^{2-n}.
\]
In particular $[n]_{\ph}\in\Z[\ph]$.  Under the chosen real embedding
\[
    \ph\mapsto \frac{1+\sqrt5}{2},
\]
the sequence $[n]_{\ph}$ is increasing and converges to $\ph^2$.
\end{proposition}

\begin{proof}
Since $q=\ph^{-1}$,
\[
    [n]_{\ph}=\frac{1-\ph^{-n}}{1-\ph^{-1}}.
\]
Using $1-\ph^{-1}=\ph^{-2}$ gives
\[
    [n]_{\ph}=\ph^2(1-\ph^{-n})=\ph^2-\ph^{2-n}.
\]
This lies in $\Z[\ph]$ because $\ph$ is a unit in $\Z[\ph]$.

Under the chosen real embedding, $\ph>1$, so $\ph^{2-n}$ is positive and decreases to $0$.  Hence $\ph^2-\ph^{2-n}$ increases to $\ph^2$.
\end{proof}

\begin{remark}[The other real embedding]
The boundedness statement is embedding-relative.  Under the conjugate embedding $\ph\mapsto\ps=1-\ph$, powers of $\ps^{-1}$ have growing magnitude and alternating sign.  Thus the sequence of conjugate weights is unbounded.  The phrase ``bounded golden weights'' always refers to the selected real place used for the analytic completion.
\end{remark}

This is the paper's first genuinely golden phenomenon.  Ordinary differentiation has monomial weights $n$, which diverge.  The golden factor derivative has monomial weights approaching $\ph^2$ at the chosen real place.

\section{Relation to Jackson \texorpdfstring{$q$}{q}-calculus}

The operator $\dd$ is formally the Jackson $q$-difference operator at
\[
    q=\ph^{-1}.
\]
Classically, the Jackson operator is often written as
\[
    \frac{f(x)-f(qx)}{x-qx}.
\]
In the present paper this display is only a comparison with standard notation: $\dd f$ is defined by the factor identity
\[
    \Del f=(\dd f)\Del x,
\]
not by a primitive pointwise quotient.  This overlap with classical $q$-calculus is intentional and should be stated openly; see Jackson's original work and standard references on quantum calculus \cite{jackson1909,kaccheung}.

The contribution here is not the invention of the $q$-difference operator.  It is the grammatical route and the golden specialization.  The contraction $x\mapsto \ph^{-1}x$ is not chosen arbitrarily; it is generated by the scalar operation
\[
    \B(\ph,x)=\ph^{-1}x.
\]
Moreover, the golden identity $1-\ph^{-1}=\ph^{-2}$ collapses the monomial weights to the bounded orbit
\[
    [n]_{\ph}=\ph^2-\ph^{2-n}.
\]
The two exponential sections below also align with the two standard $q$-exponentials.  Again, the novelty claimed here is not the pair of $q$-exponentials themselves, but their placement in the $\ph$-native layer hierarchy.

\section{The twisted product rule}

Let
\[
    \sig_{\ph}(f)(x)=f(qx).
\]
This is the contraction action on functions.

\begin{proposition}[Twisted product rule]
For $f,g\in R[x]$,
\[
    \dd(fg)=(\dd f)g+\sig_{\ph}(f)(\dd g).
\]
\end{proposition}

\begin{proof}
We compute
\begin{align*}
    \Del(fg)(x)
    &=f(x)g(x)-f(qx)g(qx)\\
    &=\bigl(f(x)-f(qx)\bigr)g(x)
      +f(qx)\bigl(g(x)-g(qx)\bigr)\\
    &=(\dd f)(x)\Del x\,g(x)+\sig_{\ph}(f)(x)(\dd g)(x)\Del x.
\end{align*}
Factoring out $\Del x$ gives the result.
\end{proof}

This is the standard $q$-Leibniz rule.  The $\ph$-native point is that the twist $\sig_{\ph}$ is the same golden contraction generated by the scalar core.

\section{Chain rule as two-stage factorization}

The ordinary chain rule does not survive unchanged.  This is not a defect; it is a structural warning about contraction calculus.

For a polynomial $f$, define the polynomial divided-difference factor $\Diff_f(A,B)$ by
\[
    f(A)-f(B)=\Diff_f(A,B)(A-B).
\]
Such a polynomial exists because $f(A)-f(B)$ is divisible by $A-B$ in $R[A,B]$: for each monomial,
\[
    A^k-B^k=(A-B)(A^{k-1}+A^{k-2}B+\cdots+B^{k-1}),
\]
and the general case follows by linearity.

\begin{proposition}[Chain factorization]
For $f,g\in R[x]$,
\[
    \dd(f\circ g)(x)=\Diff_f(g(x),g(qx))\, (\dd g)(x).
\]
\end{proposition}

\begin{proof}
We have
\begin{align*}
    \Del(f\circ g)(x)
    &= f(g(x))-f(g(qx))\\
    &= \Diff_f(g(x),g(qx))\bigl(g(x)-g(qx)\bigr)\\
    &= \Diff_f(g(x),g(qx))(\dd g)(x)\Del x.
\end{align*}
Factoring by $\Del x$ gives the formula.
\end{proof}

The classical-looking formula
\[
    \dd(f\circ g)=((\dd f)\circ g)(\dd g)
\]
would require
\[
    \Diff_f(g(x),g(qx))=(\dd f)(g(x)).
\]
But $(\dd f)(g(x))$ uses the pair $g(x)$ and $qg(x)$, whereas the actual divided difference uses the pair $g(x)$ and $g(qx)$.  These agree only when
\[
    g(qx)=qg(x).
\]

\begin{proposition}[Clean chain rule condition for polynomials]
Let $g\in R[x]$.  Then
\[
    g(qx)=qg(x)
\]
if and only if $g(x)=cx$ for some $c\in R$.
\end{proposition}

\begin{proof}
Write
\[
    g(x)=\sum_{d\ge0}a_dx^d.
\]
Then
\[
    g(qx)-qg(x)=\sum_{d\ge0}a_d(q^d-q)x^d.
\]
Since $q=\ph^{-1}$ is not a root of unity, $q^d=q$ only for $d=1$.  Therefore all coefficients except possibly $a_1$ must vanish.  Hence $g(x)=cx$.  The converse is immediate.
\end{proof}

Thus there is no general ordinary chain rule in the polynomial golden calculus.  The honest general statement is the two-stage factorization above.

\section{The reciprocal golden exponential}

A natural ordinary-looking exponential analogue is an eigenfunction of the golden derivative:
\[
    \dd E=E.
\]
Let
\[
    E(x)=\sum_{n\ge0}a_nx^n.
\]
Using
\[
    \dd x^{n+1}=[n+1]_{\ph}x^n,
\]
the eigenfunction equation yields the recurrence
\[
    a_{n+1}[n+1]_{\ph}=a_n.
\]

\begin{definition}[Formal reciprocal golden exponential]
A formal reciprocal golden exponential is a sequence $(a_n)_{n\ge0}$ satisfying
\[
    a_{0}=1,
    \qquad
    a_{n+1}[n+1]_{\ph}=a_n.
\]
\end{definition}

In the finite algebraic layer, this recurrence is the primary object.  Writing
\[
    a_n=\frac{1}{[n]_{\ph}!}
\]
requires the rate/localization layer from \emph{Quotients Without Division}.  Here
\[
    [n]_{\ph}!=[1]_{\ph}[2]_{\ph}\cdots[n]_{\ph}.
\]

After passing to a coefficient layer where the required inverses exist, define
\[
    \Eph(x)=\sum_{n=0}^{\infty}\frac{x^n}{[n]_{\ph}!}.
\]
The bounded-weight theorem immediately controls the radius of convergence.

\begin{theorem}[Radius of the reciprocal golden exponential]
Under the chosen real embedding $\ph=(1+\sqrt5)/2$, the power series $\Eph(x)$ has radius of convergence
\[
    R=\ph^2=\ph+1.
\]
In particular, $\Eph$ is not entire.
\end{theorem}

\begin{proof}
Let
\[
    a_n=\frac{1}{[n]_{\ph}!}.
\]
Then
\[
    \left|\frac{a_n}{a_{n+1}}\right|=[n+1]_{\ph}.
\]
Under the chosen real embedding,
\[
    [n+1]_{\ph}\longrightarrow \ph^2.
\]
By the ratio test, the radius of convergence is $\ph^2$.
\end{proof}

Thus the boundedness of the derivative weights and the non-entire nature of this exponential are two faces of one fact.  Ordinary calculus has monomial weights $n\to\infty$, and the analogous exponential radius is infinite.  Golden calculus has monomial weights $[n]_{\ph}\to\ph^2$, and the reciprocal exponential radius is finite.

\begin{remark}[One theorem, two faces]
The statement ``the golden derivative weights are bounded'' and the statement ``the reciprocal golden exponential is not entire'' are not independent.  The latter is the ratio-test consequence of the former.  This is the main analytic payoff of fixing the contraction ratio at $\ph^{-1}$.
\end{remark}

\begin{proposition}[Reciprocal product form]
The analytic solution of $\dd E=E$ normalized by $E(0)=1$ satisfies formally
\[
    \Eph(x)=\prod_{k\ge0}(1-\ph^{-2-k}x)^{-1}.
\]
\end{proposition}

\begin{proof}
The equation $\dd E=E$ says
\[
    E(x)-E(qx)=E(x)(x-qx).
\]
Since $x-qx=\ph^{-2}x$, this becomes
\[
    E(qx)=E(x)(1-\ph^{-2}x),
\]
or
\[
    E(x)=\frac{E(qx)}{1-\ph^{-2}x}.
\]
Iterating and using the normalization along the contraction orbit gives the displayed reciprocal product.
\end{proof}

This product form makes the radius visible: the first pole in the chosen real embedding occurs at $x=\ph^2$.

\section{The twisted product exponential}

The product rule is $\sig_{\ph}$-twisted.  This suggests a second eigen-equation, aligned with the twist:
\[
    \dd\mathcal E=\sig_{\ph}(\mathcal E).
\]
This equation is not meant to replace the previous one.  It lands in a different layer.  The untwisted equation $\dd E=E$ produces a localized reciprocal exponential.  The twisted equation produces finite polynomial approximants and then an infinite product completion.

\begin{theorem}[Twisted product exponential]
Let $q=\ph^{-1}$ and define finite products
\[
    \mathcal E_{\ph,N}(x)=\prod_{k=0}^{N}(1+\ph^{-2-k}x).
\]
Each $\mathcal E_{\ph,N}$ lies in $\Z[\ph][x]$.  In the chosen real or complex completion, the infinite product
\[
    \Etw(x)=\prod_{k\ge0}(1+\ph^{-2-k}x)
\]
converges for every $x\in\C$ and satisfies
\[
    \dd\Etw=\sig_{\ph}(\Etw).
\]
Thus $\Etw$ is entire.
\end{theorem}

\begin{proof}
The finite products have coefficients in $\Z[\ph]$ because each $\ph^{-2-k}$ is a unit of $\Z[\ph]$.

To derive the equation, write $a=\ph^{-2}=1-q$.  The twisted eigen-equation says
\[
    \mathcal E(x)-\mathcal E(qx)=\mathcal E(qx)(x-qx)=\mathcal E(qx)ax,
\]
so
\[
    \mathcal E(x)=\mathcal E(qx)(1+ax).
\]
Iterating gives
\[
    \mathcal E(x)=\prod_{k=0}^{N}(1+aq^kx)\,\mathcal E(q^{N+1}x).
\]
Since $aq^k=\ph^{-2-k}$ and $q^{N+1}x\to0$ in the chosen real or complex completion, the normalization $\mathcal E(0)=1$ yields
\[
    \mathcal E(x)=\prod_{k\ge0}(1+\ph^{-2-k}x).
\]

The product converges locally uniformly on $\C$ because for every bounded set $|x|\le M$,
\[
    \sum_{k\ge0}|\ph^{-2-k}x|
    \le M\sum_{k\ge0}\ph^{-2-k}<\infty.
\]
Hence it defines an entire function.  The functional equation above then gives $\dd\Etw=\sig_{\ph}(\Etw)$.
\end{proof}

\begin{remark}[Two standard $q$-exponentials]
In classical $q$-calculus, the reciprocal series and the product series are the two standard $q$-exponentials.  This paper does not claim to discover them.  The point is their placement in the $\ph$-native hierarchy.  The ordinary-looking eigen-equation $\dd E=E$ requires coefficient localization and has a finite analytic radius.  The twist-compatible equation $\dd\mathcal E=\sig_{\ph}(\mathcal E)$ has denominator-free finite approximants and requires completion, not coefficient localization.  Each equation is legitimate; they simply live in different layers.
\end{remark}

\begin{center}
\small
\begin{tabular}{p{0.24\linewidth}p{0.34\linewidth}p{0.26\linewidth}}
\toprule
Equation & Solution form & Layer \\
\midrule
$\dd E=E$ & $\displaystyle\prod_{k\ge0}(1-\ph^{-2-k}x)^{-1}$ & localization; radius $\ph^2$ \\
$\dd\mathcal E=\sig_{\ph}(\mathcal E)$ & $\displaystyle\prod_{k\ge0}(1+\ph^{-2-k}x)$ & polynomial/completion; entire \\
\bottomrule
\end{tabular}
\end{center}

This distinction repairs a possible discomfort in the reciprocal exponential.  The reciprocal exponential is not wrong; it is the eigenfunction that asks for localization.  The twisted product exponential is the eigenfunction whose finite approximants remain inside the native polynomial layer.

\section{Logarithm as unit-orbit index}

The product exponential also clarifies where logarithm belongs.  The tempting scalar question is to ask for a power-series logarithm or a compositional inverse.  That is not the native starting point.  The native logarithmic datum is the index of the unit orbit.

Let
\[
    U_{\ph}^{+}=\{\ph^k:k\in\Z\}
\]
be the positive golden unit group.  Define
\[
    \ord_{\ph}(\ph^k)=k.
\]
Then
\[
    \ord_{\ph}(\ph^j\ph^k)=j+k
    =\ord_{\ph}(\ph^j)+\ord_{\ph}(\ph^k).
\]
Thus logarithm's product-to-sum role is exact, integer-valued, and division-free on the unit group.  In this sense, logarithm is native on the group before it is analytic on the line.

This index is already present in the contraction orbit.  The scales in the product exponential are
\[
    \ph^{-2-k},\qquad k=0,1,2,\ldots,
\]
so the finite product
\[
    \mathcal E_{\ph,N}(x)=\prod_{k=0}^{N}(1+\ph^{-2-k}x)
\]
is built by walking along discrete logarithmic time.  The factor indexed by $k$ vanishes when
\[
    1+\ph^{-2-k}x=0,
\]
that is, at
\[
    x=-\ph^{2+k}.
\]
Hence the zeros of the product exponential lie on the negative golden unit orbit
\[
    -\ph^2,-\ph^3,-\ph^4,\ldots.
\]
The reciprocal exponential has the complementary obstruction: its poles lie on the positive orbit
\[
    \ph^2,\ph^3,\ph^4,\ldots.
\]
Thus the radius $\ph^2$ is not only a ratio-test limit.  It is the distance from the origin to the first point where the golden unit orbit produces an analytic obstruction: a pole for the reciprocal exponential and, after applying scalar logarithm, a branch point for the product exponential.

If one now applies the ordinary analytic logarithm to the product exponential, one obtains
\[
    \log \Etw(x)=\sum_{k\ge0}\log(1+\ph^{-2-k}x).
\]
This is a useful analytic coordinate, but it is no longer the primitive object.  Expanding it gives
\[
    \log \Etw(x)
    =\sum_{m\ge1}(-1)^{m+1}
      \frac{\ph^{2-2m}}{m[m]_{\ph}}x^m.
\]
The denominators $m[m]_{\ph}$ record the cost of interpolating the integer unit index off the unit group.  The integer factor $m$ comes from ordinary analytic interpolation, while the golden factor $[m]_{\ph}$ comes from summing over the contraction orbit.  The regulator $\log\ph$ is the archimedean step-size of this interpolation.

This is not a new logarithm in the classical sense.  The unit index, regulators, valuations, and $q$-products are classical objects.  The point here is their placement in the $\ph$-native hierarchy: logarithm is exact and denominator-free as an orbit index on the unit group; scalar analytic logarithm appears only after that index is extended away from the discrete orbit.  In the elliptic regime of the phase paper, the analogous index becomes cyclic, such as the $\Z/10\Z$ index of the golden phase; we leave that finite-order logarithmic picture for later work.

\section{Integration as rate localization}

A $\ph$-antiderivative of $x^n$ asks for a polynomial or rate object $F$ satisfying
\[
    \dd F=x^n.
\]
If we try
\[
    F=a x^{n+1},
\]
then
\[
    \dd(a x^{n+1})=a[n+1]_{\ph}x^n.
\]
So we need
\[
    a[n+1]_{\ph}=1.
\]
Thus integration requires localization at the golden weights.

Let
\[
    S_{\ph}=\langle [1]_{\ph},[2]_{\ph},[3]_{\ph},\ldots\rangle
\]
be the multiplicative set generated by the golden weights.  Then monomial integration naturally lives over
\[
    S_{\ph}^{-1}R.
\]
Before evaluating the coefficient as a scalar, we can write the primitive as a rate object:
\[
    \int_{\ph}x^n\,d_{\ph}x
    =[x^{n+1}:[n+1]_{\ph}]_{S_{\ph}}.
\]

The weights themselves have a simple unit classification.

\begin{proposition}[Units among the golden weights]
For $n\ge1$, the weight $[n]_{\ph}$ is a unit in $\Z[\ph]$ if and only if $n=1$ or $n=2$.
\end{proposition}

\begin{proof}
Using
\[
    [n]_{\ph}=\ph^{2-n}(\ph^n-1),
\]
and the fact that $\ph$ is a unit, $[n]_{\ph}$ is a unit if and only if $\ph^n-1$ is a unit.  Its norm is
\begin{align*}
    \Nm(\ph^n-1)
    &=(\ph^n-1)(\ps^n-1)\\
    &=(\ph\ps)^n-(\ph^n+\ps^n)+1\\
    &=(-1)^n-L_n+1,
\end{align*}
where
\[
    L_n=\ph^n+\ps^n
\]
is the Lucas number.  Therefore
\[
    \Nm([n]_{\ph})=(-1)^n\bigl((-1)^n-L_n+1\bigr).
\]
Equivalently,
\[
    \Nm([n]_{\ph})=
    \begin{cases}
        L_n, & n\text{ odd},\\
        2-L_n, & n\text{ even}.
    \end{cases}
\]
For $n=1$, this gives $\Nm([1]_{\ph})=1$.  For $n=2$, this gives $\Nm([2]_{\ph})=-1$, and explicitly
\[
    [2]_{\ph}=1+\ph^{-1}=1+(\ph-1)=\ph.
\]
Thus $[1]_{\ph}=1$ and $[2]_{\ph}=\ph$ are units.  For odd $n\ge3$, $L_n\ge4$, so the norm is not $\pm1$.  For even $n\ge4$, $L_n\ge7$, so $2-L_n\le -5$, again not $\pm1$.  Hence no other golden weights are units.
\end{proof}

Thus the integration obstruction is generic, not accidental.  Apart from the first two weights, antiderivatives require a genuine rate/localization layer.

\section{Fixed scale, not infinitesimal limit}

This calculus is not a replacement for ordinary infinitesimal differentiation.  It is a self-similar finite-scale calculus attached to the native contraction
\[
    x\mapsto \ph^{-1}x.
\]

The ordinary derivative appears in a different regime, where one studies a family of $q$-difference operators and lets $q\to1$.  Here $q$ is fixed at $\ph^{-1}$.  The fixed scale explains both the algebraic factor projection and the analytic behavior of the exponentials.  Because the scale does not shrink to $1$, the weights do not grow like $n$; they approach $\ph^2$ at the selected real place.  Consequently, the reciprocal exponential has finite radius.  Meanwhile, the twist-compatible product exponential remains denominator-free at finite depth and becomes entire only after completion.  The same unit orbit controls the logarithmic boundary: poles occur on the positive orbit for the reciprocal exponential, and zeros/branch points occur on the negative orbit for the product exponential and its scalar logarithm.

This is honesty, but it is also structure:
\begin{quote}
ordinary calculus is infinitesimal; golden calculus is self-similar.
\end{quote}

\section{Comparison with classical calculus}

\begin{center}
\small
\begin{tabular}{p{0.42\linewidth}p{0.46\linewidth}}
\toprule
Classical calculus & Golden-native calculus \\
\midrule
Additive displacement $x+h$ & Multiplicative contraction $qx$, $q=\ph^{-1}$ \\
Limit $h\to0$ & Fixed self-similar scale \\
Derivative quotient & Displacement factor projection \\
Monomial weight $n$ & Bounded real-place weight $[n]_{\ph}\to\ph^2$ \\
Ordinary product rule & Twisted $q$-Leibniz rule \\
Ordinary chain rule & Two-stage factorization; clean rule only for $g(x)=cx$ \\
Exponential $e^x$ entire & Reciprocal exponential radius $\ph^2$; product exponential entire \\
Logarithm as scalar function & Unit-orbit index on $\{\ph^k\}$; analytic interpolation off the group \\
Integral coefficient $1/(n+1)$ & Rate localization at $[n+1]_{\ph}$ \\
\bottomrule
\end{tabular}
\end{center}

The table should not be read as a claim that golden calculus supersedes ordinary calculus.  It is a different calculus, adapted to a different primitive displacement.

\section{Open questions}

\begin{question}[Golden Taylor theory]
Is there a useful Taylor theory based on iterated golden contractions
\[
    x,
    \ qx,
    \ q^2x,
    \ldots?
\]
The answer is not automatic: diagonal Taylor expansion and contraction-orbit interpolation are different problems.
\end{question}

\begin{question}[Function spaces]
Which function spaces beyond polynomials admit a natural golden derivative?  Power series inside the disk of radius $\ph^2$ are one candidate for the reciprocal exponential, while the product exponential suggests entire functions built from contraction-orbit products.  The correct analytic category remains to be determined.
\end{question}

\begin{question}[Unit-index logarithms]
Can the unit-index interpretation of logarithm be connected systematically to the analytic logarithm through archimedean completion and to valuations through non-archimedean completions?  In particular, can the radius $\ph^2$ of the reciprocal/product logarithmic obstructions be characterized entirely as the first obstruction along the golden unit orbit?
\end{question}

\begin{question}[Weight arithmetic and the Lucas semigroup]
How do the elements
\[
    [n]_{\ph}=\ph^2-\ph^{2-n}
\]
factor in $\Z[\ph]$?  The unit classification above shows that only $n=1,2$ are units, and the proof shows that their norms are governed by Lucas numbers:
\[
    \Nm([n]_{\ph})=
    \begin{cases}
        L_n, & n\text{ odd},\\
        2-L_n, & n\text{ even}.
    \end{cases}
\]
Thus the integration localization is controlled by the multiplicative semigroup generated by the golden weights, whose norms lie in a Lucas-number semigroup.  How does this ``Lucas semigroup'' relate to Lucas divisibility, ranks of apparition, and the splitting or ramification of rational primes in $\Q(\sqrt5)$?  The prime factorization behavior of these weights controls the arithmetic of golden integration.
\end{question}

\begin{question}[Two exponential layers]
Can the reciprocal and product exponentials be characterized intrinsically as representing two universal properties: localization for $\dd E=E$, and completion for $\dd\mathcal E=\sig_{\ph}(\mathcal E)$?  How far does this layer distinction persist for other $q$-difference calculi?
\end{question}

\begin{question}[Golden differential equations]
Which equations are naturally solved by the non-entire reciprocal exponential $\Eph$, and which are naturally solved by the entire product exponential $\Etw$?  How do their analytic properties differ from ordinary exponential equations?
\end{question}

\begin{question}[Deformation to ordinary calculus]
Can the golden operator be placed inside a family of factor projections depending on $q$, with ordinary calculus recovered as $q\to1$?  How does the bounded-weight phenomenon disappear in that limit?
\end{question}

\begin{question}[Chain-rule geometry]
For polynomial maps, the clean chain rule holds only for $g(x)=cx$.  Is there a useful geometric class of non-polynomial maps commuting with golden contraction in an appropriate analytic category?
\end{question}

\section{Conclusion: the derivative as a factor shadow}

Classically, calculus is introduced through a quotient.  In the polynomial setting, however, the quotient exists because a divisibility relation is already present:
\[
    f(x)-f(qx)
\]
is divisible by
\[
    x-qx.
\]
The golden derivative begins from that factor relation.

The golden leaf contributes more than a suggestive basis.  It supplies a native contraction
\[
    x\mapsto\ph^{-1}x
\]
and forces the derivative weights into the closed form
\[
    [n]_{\ph}=\ph^2-\ph^{2-n}.
\]
At the chosen real place, these weights are bounded and converge to $\ph^2$.  The associated reciprocal exponential therefore has radius $\ph^2$ and is not entire.

The twisted product rule reveals the complementary exponential.  The equation $\dd E=E$ lands in the localization layer and produces a reciprocal product.  The equation $\dd\mathcal E=\sig_{\ph}(\mathcal E)$ lands first in the polynomial/completion layer and produces denominator-free finite products.  These are the two standard $q$-exponentials at the golden value $q=\ph^{-1}$, but the $\ph$-native hierarchy explains why they belong to different layers.

The logarithmic story collapses to the same orbit.  On the unit group, logarithm is the exact index $\ord_{\ph}(\ph^k)=k$.  Off the group, scalar analytic logarithm is an interpolation of that index, and the denominators in its series record the cost of leaving the discrete orbit.  The radius $\ph^2$ is the nearest point of the same golden unit orbit: a pole for the reciprocal exponential, and a zero or logarithmic branch point for the product exponential.

Thus the fixed scale is not a defect to be hidden.  It is the source of the calculus.  The derivative is not born as a quotient; it is the shadow of a factor, and logarithm is native first as an index before it is analytic as a function.

\begin{thebibliography}{99}

\bibitem{santosGoldenShadow}
V. Santos.
\newblock A golden shadow algebra: trace completion from a \(\varphi\)-leaf scalar grammar.
\newblock Companion manuscript, 2026.

\bibitem{santosQuotients}
V. Santos.
\newblock Quotients without division: unit actions, rate localizations, and factor projections in a \(\varphi\)-native algebra.
\newblock Companion manuscript, 2026.

\bibitem{santosShadowTrig}
V. Santos.
\newblock Shadow trigonometry: norm-one orbits, golden traces, and cyclotomic phase.
\newblock Companion manuscript, 2026.

\bibitem{santosSignatures}
V. Santos.
\newblock Signatures before roots: shadow lifts, norm forms, and radical spreads.
\newblock Companion manuscript, 2026.

\bibitem{jackson1909}
F. H. Jackson.
\newblock On $q$-functions and a certain difference operator.
\newblock \emph{Transactions of the Royal Society of Edinburgh}, 46:253--281, 1909.

\bibitem{kaccheung}
V. Kac and P. Cheung.
\newblock \emph{Quantum Calculus}.
\newblock Universitext. Springer, 2002.

\bibitem{odrzywolek}
A. Odrzywo\l{}ek.
\newblock All elementary functions from a single binary operator.
\newblock arXiv:2603.21852, 2026.

\end{thebibliography}

\end{document}
